\documentclass[11pt, a4paper]{article}
\usepackage[UTF8]{ctex} % 支持中文
\usepackage{amsmath}    % 数学公式
\usepackage{amsfonts}
\usepackage{geometry}
\geometry{left=2.5cm,right=2.5cm,top=2.5cm,bottom=2.5cm}

\begin{document}


\begin{enumerate}
    % 第一题：判断题
    \item \textbf{【判断题】（难度：低）} \\
    在进行统计建模时，如果因变量 $Y$ 是定性数据（例如：消费者是否购买某产品、病人的诊断结果为“得病”或“未得病”），我们通常可以直接使用经典的线性回归模型进行建模，而不需要做任何特殊处理。 \hfill （\quad）

    \vspace{2em}

    % 第二题：选择题
    \item \textbf{【选择题】（难度：中）} \\
    在评估逻辑回归或其他二分类模型的分类能力时，常会使用 ROC 曲线下的面积（AUC 值）。关于 AUC，下列说法中正确的是？
    \begin{enumerate}
        \item AUC 的取值大小会随着分类阈值（Threshold）的改变而发生变化。
        \item AUC 值越接近 0.5，说明该模型的分类预测效果越理想。
        \item AUC 衡量的是模型在所有可能的分类阈值下的平均表现，其取值与阈值的选取无关。
        \item 当 AUC 等于 1 时，说明模型对所有样本的预测结果与实际情况完全相反。
    \end{enumerate}

    \vspace{2em}

    % 第三题：填空题
    \item \textbf{【填空题】（难度：高）} \\
    在生存分析中，由于研究结束、失访或个体死于其他原因，导致我们无法观测到个体确切的“失效”或“死亡”时间，这种现象在统计学上被称为\underline{\hspace{4em}}。在生存回归模型中，\underline{\hspace{4em}}模型（比例风险模型）是对“风险”进行建模，其估计出来的系数符号通常与 AFT（加速失效时间）模型估计出的系数符号\underline{\hspace{4em}}。

\end{enumerate}

\vspace{4em}

% 答案部分（可选）
\begin{hr}
\end{hr}
\textbf{参考答案及解析：}
\begin{itemize}
    \item \textbf{1. 错误。} 解析：定性数据应选用逻辑回归、决策树等分类模型，而非线性回归。
    \item \textbf{2. C。} 解析：AUC 反映模型整体区分能力，与阈值无关；0.5 表示随机水平。
    \item \textbf{3. 数据删失（Censoring）；Cox；相反。} 解析：Cox 针对风险建模，AFT 针对时间建模，风险增加意味着存活时间缩短，故系数符号通常相反。
\end{itemize}

\end{document}
